\section{Joint Distribution}

Recall that, if \(X\) is a discrete random variable then the distribution of \(X\) is \(\para{\Prob\para{X = x}}_x\).

\subsection{Joint Distribution}

\begin{definition}[Joint Distribution for Discrete Variables]
    \label{def:joint-distribution-discrete}%
    Let \(X_1, \ldots, X_n\) be discrete random variables. Their \emph{joint distribution} is defined to be
    \[
        \Prob\para{X_1 = x_1, \ldots, X_n = x_n}
    \]
    for all possible values of \(x_1, x_2, \ldots, x_n\).
\end{definition}

Note that we could recover the individual distribution of a single random variable from this as well. Indeed, we have
\[
    \Prob\para{X_1 = x_1} = \sum_{x_2, \ldots, x_n \in \RR} \Prob\para{X_1 = x_1, \ldots, X_n = x_n}
\]
where we note that this sum is countable since \(x_i\)s only take countably many values. We call this the \emph{marginal distribution} of \(X_1\).

\begin{definition}[Conditional Distribution for Discrete Variables]
    \label{def:conditional-distribution-discrete}%
    Let \(X\) and \(Y\) be two discrete random variables. The \emph{conditional distribution} of \(X\), given \(Y = y\) (where \(y\) is fixed) is defined to be
    \[
        \Prob\para{\Conditional{X = x}{Y = y}} = \frac{\Prob\para{X = x, Y = y}}{\Prob\para{Y = y}}
    \]
    for all possible values of \(x\).
\end{definition}

By \autoref{thm:law-total-probability} (Law of Total Probability), we have
\[
    \Prob\para{X = x} = \sum_{y \in \RR} \Prob\para{\Conditional{X = x}{Y = y}} \cdot \Prob\para{Y = y}
\]
noting that this sum is similarly countable.

\subsection{Sum of Random Variables}

\begin{example}[Distribution of Sum of Random Variables]
    \label{ex:distribution-sum-discrete}%
    Let \(X\) and \(Y\) be two independent random variables. We = want to understand the distribution of \(X + Y\). We have
    \begin{align*}
        \Prob\para{X + Y = z} &= \sum_{y \in \RR} \Prob\para{X + Y = z, Y = y} \\
        &= \sum_{y \in \RR} \Prob\para{X = z - y, Y = y} \\
        &= \sum_{y \in \RR} \Prob\para{X = z - y} \cdot \Prob\para{Y = y}.
    \end{align*}
\end{example}

\begin{definition}[Convolution of Distribution for Discrete Variables]
    \label{def:convolution-distribution-discrete}%
    Let \(X, Y\) be two discrete random variables. Then the distribution of \(X + Y\), given by
    \[
        \sum_{y \in \RR} \Prob\para{X = z - y} \cdot \Prob\para{Y = y}
    \]
    is known as the \emph{convolution} of the two random variables.
\end{definition}

\begin{example}[Sum of Poisson Random Variables]
    Let \(X \sim \Poisson\para{\lambda}\), \(Y \sim \Poisson\para{\mu}\), \(\lambda, \mu > 0\) and \(X \indp Y\). We may show that the sum of two independent Poisson random variables is another Poisson random variable.

    \begin{align*}
        \Prob\para{X + Y = n} &= \sum_{k = 0}^{n} \Prob\para{X = n - k} \cdot \Prob\para{Y = k} \\
        &= \sum_{k = 0}^{m} e^{-\lambda} \cdot \frac{\lambda^{n - k}}{\para{n - k}!} \cdot e^{-\mu} \cdot \frac{\mu^k}{k!} \\
        &= e^{-\para{\lambda + \mu}}{n!} \cdot \sum_{k = 0}^{n} \binom{n}{k} \cdot \lambda^{n - k} \cdot \mu^k \\
        &= e^{-\para{\lambda + \mu}} \cdot \frac{\para{\lambda + \mu}^n}{n!}.
    \end{align*}

    Therefore, we have
    \[
        \Prob\para{X + Y = n} = e^{-\para{\lambda + \mu}} \cdot \frac{\para{\lambda + \mu}^n}{n!}
    \]
    and hence
    \[
        X + Y \sim \Poisson \para{\lambda + \mu}.
    \]

    Indeed, if \(\lambda = \mu\), let \(X_n, Y_n \sim \Binomial\para{n, \frac{\lambda}{n}}\) and \(X_n \indp Y_n\). We know a Poisson distribution is a limit of binomial random variables, and indeed
    \[
        X_n + Y_n \sim \Binomial\para{2n, \frac{\lambda}{n}}.
    \]
\end{example}

\subsection{Conditional Expectation}

\begin{definition}[Conditonal Expectation upon an Event]
    \label{def:conditional-expectation-upon-event}%
    The \emph{conditional expectation} of a random variable \(X\) with respect to an event \(B\), where \(\Prob\para{B} > 0\), is given by
    \[
        \Expt\brac{\Conditional{X}{B}} = \frac{\Expt\brac{X \cdot \indicator\para{B}}}{\Prob\para{B}}.
    \]
\end{definition}

Note that with \(X = \indicator\para{A}\), we do recover consistent definitions with conditional probability of events.

\[
    \Expt\brac{\Conditional{\indicator\para{A}}{B}} = \frac{\Expt\brac{\indicator\para{A} \cdot \indicator\para{B}}}{\Prob\para{B}} = \frac{\Expt\brac{\indicator\para{A \cap B}}}{\Prob\para{B}} = \frac{\Prob\para{A \cap B}}{\Prob\para{B}} = \Prob\para{\Conditional{A}{B}}.
\]

\begin{theorem}[Law of Total Expectation]
    \label{thm:law-of-total-expectation}%
    Let \(X\) be a random variable, \(\para{\Omega_i}\) be a partition of \(\Omega\) into disjoint sets, with \(\Prob\para{\Omega_i} > 0\) for all \(i = 1, 2, \ldots ,n\). Then we have
    \[
        \Expt\brac{X} = \sum_{i = 1}^{n} \Expt\brac{\Conditional{X}{\Omega_n}} \cdot \Prob\para{\Omega_n}.
    \]
\end{theorem}

\begin{proof}
    We have
    \begin{align*}
        \Expt\brac{X} &= \Expt\brac{X \cdot \indicator\para{\Omega}} \\
        &= \Expt\brac{X \cdot \sum_{i = 1}^{n} \indicator \para{\Omega_i}} \\
        &= \sum_{i = 1}^{n} \Expt\brac{X \cdot \indicator \para{\Omega_i}} \\
        &= \sum_{i = 1}^{N} \Expt\brac{\Conditional{X}{\Omega_n}} \cdot \Prob\para{\Omega_n}
    \end{align*}
\end{proof}

The conditional expectation of \(X\) given \(Y = y\) is defined similarly, since \(Y = y\) is an event:
\[
    \Expt\brac{\Conditional{X}{Y = y}} = \frac{\Expt\brac{X \cdot \indicator\para{Y = y}}}{\Prob\para{Y = y}}.
\]

\begin{proposition}[Conditional Expectation upon Random Variable Event]
    \label{prop:conditional-expectation-random-variable}%
    We have
    \[
        \Expt\brac{\Conditional{X}{Y = y}} = \sum_{x \in \RR} x \cdot \Prob\para{\Conditional{X = x}{Y = y}}.
    \]
\end{proposition}

\begin{proof}
    We have
    \begin{align*}
        \Expt\brac{\Conditional{X}{Y = y}} &= \frac{\Expt\brac{X \cdot \indicator\para{Y = y}}}{\Prob\para{Y = y}} \\
        &= \sum_{x \in \RR} x \cdot \frac{\Prob\para{X = x, Y = y}}{\Prob\para{Y = y}} \\
        &= \sum_{x \in \RR} x \cdot \Prob\para{\Conditional{X = x}{Y = y}}.
    \end{align*}
\end{proof}

This shows that the conditional expectation is consistent with the expectation of the conditional random variable.

\begin{definition}[Conditional Expectation Distribution]
    \label{def:conditional-expectation-distribution}%
    We define \(\Expt\brac{\Conditional{X}{Y}}\) as a random variable as
    \[
        \Expt\brac{\Conditional{X}{Y}} = g\para{Y} = \sum_{y \in \RR} g\para{y} \cdot \indicator\para{Y = y},
    \]
    where
    \[
        g\para{y} = \Expt\brac{\Conditional{X}{Y = y}}.
    \]
\end{definition}

This notation may seem confusing at the start. For example, if \(g\para{y} = y^2\), we have
\[
    \Expt\brac{\Conditional{X}{Y}} = g\para{Y} = Y^2.
\]

Intuitively, this is the \textbf{best approximation} of \(X\) using the random variable \(Y\). In the two extreme cases, if \(X = f\para{Y}\), then we have
\[
    \Expt\brac{\Conditional{X}{Y}} = f\para{Y},
\]
and if \(X \indp Y\), we have
\[
    \Expt\brac{\Conditional{X}{Y}} = \Expt\brac{X}.
\]

As a function, we may define it for \(\omega \in \Omega\),
\[
    \Expt\brac{\Conditional{X}{Y}}\para{\omega} = \Expt\brac{\Conditional{X}{Y = Y\para{\omega}}}.
\]

\begin{example}[Conditional Expectation for Tossing \(p\)-Coin]
    \label{ex:conditional-expectatio-tossing-p-coin}%
    Toss a \(p\)-coin independently for \(n\) times. Let \(X_i\) be the indicator that the \(i\)th toss is a head, for \(i = 1, 2, \ldots, n\), and let \(Y_n = X_1 + X_2 + \cdots + X_n\).

    We set \(g\para{y} = \Expt\brac{\Conditional{X_1}{Y_n = y}}\). Note that
    \begin{align*}
        g\para{y} &= \Expt\brac{\Conditional{X_1}{Y_n = y}} \\
        &= \frac{\Expt\brac{X_1 \cdot \indicator\para{Y_n = y}}}{\Prob\para{Y_n = y}} \\
        &= \frac{\Prob\para{X_1 = 1, Y_n = y}}{\Prob\para{Y_n = y}} \\
        &= \frac{p \cdot \binom{n - 1}{y - 1} \cdot p^{y - 1} \cdot \para{1 - p}^{n - y}}{\binom{n}{y} \cdot p^y \cdot \para{1 - p}^{n - y}} \\
        &= \frac{y}{n}.
    \end{align*}

    Hence,
    \[
        \Expt\brac{\Conditional{X_1}{Y_n}} = \frac{Y_n}{n}.
    \]

    Indeed, this makes sense, since one would probably expect each \(X_i\) to be \(\frac{Y}{n}\) `on average'.
\end{example}

\begin{proposition}[Properties of Conditional Expectation Distribution]
    \label{prop:conditional-expectation-distribution-property}%
    For any \(c \in \RR\), we have
    \begin{itemize}
        \item \(\Expt\brac{\Conditional{cX}{Y}} = c \Expt\brac{\Conditional{X}{Y}}\).
        \item \(\Expt\brac{\Conditional{c + X}{Y}} = c + \Expt\brac{\Conditional{X}{Y}}\) and \(\Expt\brac{\Conditional{c}{Y}} = c\).
        \item If \(X_1, X_2, \ldots, X_n\) and \(Y\) are random variables, then we have
        \[
            \Expt\brac{\Conditional{\sum_{i = 1}^{n} X_i}{Y}} = \sum_{i = 1}^{n} \Expt\brac{\Conditional{X_i}{Y}}.
        \]
    \end{itemize}
\end{proposition}

\begin{proposition}[Expectation of Conditional Expectation Distribution]
    \label{prop:expectation-conditional-expectation-distribution}%
    We have
    \[
        \Expt\brac{\Expt\brac{\Conditional{X}{Y}}} = \Expt\brac{X}.
    \]
\end{proposition}

\begin{proof}
    Recall that
    \[
        \Expt\brac{\Conditional{X}{Y}} = \sum_{y \in \RR} \indicator\para{Y = y} \Expt\brac{\Conditional{X}{Y = y}}.
    \]

    Taking the expectations on both sides, we have
    \[
        \Expt\brac{\Expt\brac{\Conditional{X}{Y}}} = \sum_{y \in \RR} \Prob\para{Y = y} \cdot \Expt\brac{\Conditional{X}{Y = y}} = \Expt\brac{X}.
    \]
\end{proof}

\begin{proposition}[Conditional Expectation Distribution of Independent Variables]
    \label{prop:conditional-expectation-distribution-independent}%
    Let \(X\) and \(Y\) be independent. Then
    \[
        \Expt\brac{\Conditional{X}{Y}} = \Expt\brac{X}.
    \]
\end{proposition}

\begin{proof}
    Recall that by definition,
    \begin{align*}
        \Expt\brac{\Conditional{X}{Y = y}} &= \sum_{x \in \RR} x \cdot \Prob\para{\Conditional{X = x}{Y = y}} \\
        &= \sum_{x \in \RR} x \cdot \Prob\para{X = x} \\
        &= \Expt\brac{X}.
    \end{align*}

    Hence,
    \[
        \Expt\brac{\Conditional{X}{Y}} = \Expt\brac{X} \sum_{y \in \RR} \indicator\para{Y = y} = \Expt\brac{X}
    \]
    as desired.
\end{proof}

\begin{proposition}[Conditional Expectation upon Independent Conditions]
    \label{prop:conditional-expectation-independent-conditions}%
    Let \(Y\) and \(Z\) be two independent random variables. Then
    \[
        \Expt\brac{\Conditional{\Expt\brac{\Conditional{X}{Y}}}{Z}} = \Expt\brac{X}.
    \]
\end{proposition}

\begin{proof}
    Set \(g\para{Y} = \Expt\brac{\Conditional{X}{Y}}\), while \(g\para{y} = \Expt\brac{\Conditional{X}{Y = y}}\).

    First, we claim that \(g\para{Y}\) is independent of \(Z\). We want to show that for all \(w, z\), we have
    \[
        \Prob\para{g\para{Y} = w, Z = z} = \Prob\para{g\para{Y} = w} \cdot \Prob\para{Z = z}.
    \]

    We have
    \begin{align*}
        \Prob\para{g\para{Y} = w, Z = z} &= \sum_{y \in g\inverse\para{\brce{w}}} \Prob\para{Y = y, Z = z} \\
        &= \sum_{y \in g\inverse\para{\brce{w}}} \Prob\para{Y = y} \cdot \Prob\para{Z = z} \\
        &= \Prob\para{Z = z} \cdot \sum_{y \in g\inverse\para{\brce{w}}} \Prob\para{Y = y} \\
        &= \Prob\para{Z = z} \cdot \Prob\para{g\para{Y} = w}.
    \end{align*}

    Therefore, \(g\para{Y} \indp Z\).

    Therefore, by \autoref{prop:conditional-expectation-distribution-independent} we have
    \[
        \Expt\brac{\Conditional{\Expt\brac{\Conditional{X}{Y}}}{Z}} = \Expt\brac{\Conditional{g\para{Y}}{Z}} = \Expt\brac{g\para{Y}},
    \]
    and we have from \autoref{prop:expectation-conditional-expectation-distribution},
    \[
        \Expt\brac{g\para{Y}} = \Expt\brac{\Expt\brac{\Conditional{X}{Y}}} = \Expt\brac{X}.
    \]
\end{proof}

\begin{proposition}[Extraction of Function of Condition]
    \label{prop:extraction-function-condition}%
    For any function \(\functiondomain{h}{\RR}{\RR}\), we have
    \[
        \Expt\brac{\Conditional{h\para{Y} \cdot X}{Y}} = h\para{Y} \cdot \Expt\brac{\Conditional{X}{Y}}.
    \]
\end{proposition}

\begin{proof}
    We have
    \begin{align*}
        \Expt\brac{\Conditional{h\para{Y} \cdot X}{Y}} &= \sum_{y \in \RR} \indicator\para{Y = y} \cdot \Expt\brac{\Conditional{h\para{Y} \cdot X}{Y = y}} \\
        &= \sum_{y \in \RR} \indicator\para{Y = y} \cdot h\para{Y} \cdot \Expt\brac{\Conditional{X}{Y = y}} \\
        &= h\para{Y} \cdot \sum_{y \in \RR} \indicator\para{Y = y} \cdot \Expt\brac{\Conditional{X}{Y = y}} \\
        &= h\para{Y} \cdot \Expt\brac{\Conditional{X}{Y}}.
    \end{align*}
\end{proof}

We have two corollaries due to the properties above.

\begin{corollary}[Conditional Expectation upon Same Condition]
    \label{cor:conditional-expecation-upon-same-condition}%
    We have
    \[
        \Expt\brac{\Conditional{\Expt\brac{\Conditional{X}{Y}}}{Y}} = \Expt\brac{\Conditional{X}{Y}}.
    \]
\end{corollary}

\begin{corollary}[Conditional Expectation upon Itself]
    \label{cor:conditional-expectation-upon-itself}%
    We have
    \[
        \Expt\brac{\Conditional{X}{X}} = X.
    \]    
\end{corollary}

\begin{example}[Conditional Expectation for Tossing \(p\)-Coin, Revisited]
    \label{ex:conditional-expectatio-tossing-p-coin-revisited}%
    We revisit \autoref{ex:conditional-expectatio-tossing-p-coin}. We toss a \(p\)-coin \(n\) times independently. Let \(X_i = \indicator\para{i\text{-th toss is H}}\) for \(i = 1, 2, \ldots, n\), and let
    \[
        Y_n = X_1 + \cdots + X_n.
    \]

    By symmetry, for all \(i = 1, 2, \ldots, n\), we have
    \[
        \Expt\brac{\Conditional{X_i}{Y_n}} = \Expt\brac{\Conditional{X_1}{Y_n}}
    \]

    Note that
    \[
        \Expt\brac{\Conditional{\sum_{i = 1}^{n} X_i}{Y_n}} = \Expt\brac{\Conditional{Y_n}{Y_n}} = Y_n,
    \]
    while using linearity gives
    \[
        \Expt\brac{\Conditional{\sum_{i = 1}^{n} X_i}{Y_n}} = \sum_{i = 1}^{n} \Expt\brac{\Conditional{X_i}{Y_n}} = n \cdot \Expt\brac{\Conditional{X_1}{Y_n}}.
    \]

    Hence,
    \[
        \Expt\brac{\Conditional{X_1}{Y_n}} = \frac{Y_n}{n}.
    \]
\end{example}